% Input fragment; independently audited construction.
\begin{proposition}[An elliptic triple core]
\label{prop:elliptic-triple-core}
Let $\ell\ge53$ be prime and put
\[
 B=\ell-1,\quad N_0=\frac{B(B-4)}3,\quad A=B-4,\quad T=B-3,
\]
\[
 n=\frac{T^2+1}{2},\qquad
 t=n-N_0=\frac{B^2-10B+30}{6}.
\]
For infinitely many primes $p$, there is an $n$-point domain in
$\mathbb F_p$ and an affine received line $F+zG$, over $\mathbb F_p$,
with
\[
 \operatorname{agr}_3(F)=\operatorname{agr}_3(G)
 =\operatorname{CA}_3(F,G)=A.
\]
Exactly $M=Bt+1$ nonzero labels have agreement at least $T$, and each
qualifying list consists of one quadratic. The threshold satisfies
\[
 na_1(3/n)<T<\sqrt{2n},\qquad T^2=2n-1.
\]
As $\ell$ tends to infinity,
\[
 \frac{M}{n^{3/2}}\longrightarrow\frac{\sqrt2}{3},
\]
a factor $4/3$ above the leading constant of
Theorem~\ref{thm:prime-quadratic-intersection-line}.
\end{proposition}
\begin{proof}
Fix a nonsingular elliptic curve over $\mathbb Q$, embedded as a plane
cubic with identity $O$, and choose a point of order $\ell$ over the
algebraic closure. Let $H$ be its cyclic subgroup. Use the $B$ points
$H\setminus\{O\}$ as parameter points, retaining only lines through
three distinct such points whose sum is $O$.

For a fixed nonzero $R\in H$, a second point $S$ ranges over
$H\setminus\{O,R\}$. Exclude $S=-R$, $S=-2R$, and $S=-R/2$, which
would make the third point zero or repeat one of the first two. These
three exclusions are distinct because $\ell>3$. The remaining
$\ell-5=B-4$ choices are paired by $S\leftrightarrow-R-S$.
Thus each parameter point lies on $(B-4)/2$ retained lines, and there
are $B(B-4)/6$ such lines. Each contains exactly three parameter points,
by intersection with the nonsingular cubic.

Choose a projective line at infinity avoiding all parameter points and
all intersections of two retained lines. Choose horizontal and vertical
directions avoiding the retained lines as well. A rational projective
coordinate change satisfying these finitely many nonvanishing
conditions exists by Zariski density. In the resulting affine chart,
write the parameter points as $(a_i,b_i)$ and set
\[
 P_i(U)=a_iU^2+b_i\qquad(1\le i\le B).
\]
The retained lines now have distinct finite nonzero slopes. A retained
line $b=ma+c$ supplies two evaluation coordinates
$U=\pm\sqrt{-m}$ with received value $c$. Adjoin all these square roots
to the number field of the construction. The coordinates are distinct,
and each has exactly three matching bank polynomials.

This core therefore has $N_0$ coordinates and exactly $A$ matches for
each $P_i$. Any other quadratic has at most $\lfloor2B/3\rfloor$ core
matches: each match contributes three roots to its differences with
the bank, while these $B$ nonzero differences contribute at most $2B$
roots in total.

Pass to a finite Galois number field containing the construction.
Discard the finitely many primes at which a needed denominator,
distinctness condition, or nonzero coefficient degenerates, including
characteristic two. At every remaining completely split prime all
coordinates and incidences reduce to the prime field. There are
infinitely many such primes; choose one also exceeding $5B^4$.

Set $g=0$ on the core, and denote its received word by $f$. Append $t$
fresh nonzero coordinates $x$, setting $f(x)=x^4$ and $g(x)=x^3$.
Choose them greedily so all labels
\[
 \lambda_{i,x}=\frac{P_i(x)-x^4}{x^3}
\]
are distinct and outside $\{0,1\}$. In addition to zero and previous
coordinates, avoid every pairwise bank intersection, including those
not retained in the core; there are at most $B(B-1)$ such coordinates.
After $j<t<B^2$ choices, the number of forbidden points is at most
\[
 B(B-1)+1+j+8B+4jB^2<5B^4<p.
\]
The last two terms exclude labels zero and one and all earlier labels,
using degree-four root bounds. Within a fresh coordinate, label
collisions are already excluded by avoiding bank intersections.

For every $\lambda$, a nonbank quadratic matches $f+\lambda g$ at most
\[
 \lfloor2B/3\rfloor+4\le B-4=A
\]
times. Each bank member gains exactly one fresh match at each of its
$t$ designated labels and none at any other label. Thus precisely $Bt$
labels have agreement $T$, with singleton threshold lists, and every
other label has nearest agreement $A$.
The zero polynomial has $N_0$ matches with $g$; any nonzero quadratic
has at most two core and three fresh matches with $g$. Since $N_0>T>5$,
the list for $g$ at threshold $T$ consists only of zero. Also
$\operatorname{CA}_3(f,g)=A$: the pair $(P_i,0)$ attains it, and a
nonzero explanation of $g$ has at most five matches.

As in Theorem~\ref{thm:prime-quadratic-intersection-line}, set
$F=f$ and $G=f+g$. Their individual agreements are $A$, and common
agreement is preserved. The $Bt$ finite exceptions transform by
$z=\lambda/(1-\lambda)$; the additional exception $z=-1$ gives $-g$,
whose only threshold witness is zero. This proves the exact count and
all singleton assertions.

Finally, $T^2=2n-1$. Since $T\ge49$, the low-rate estimate gives
\[
 na_1(3/n)\le\sqrt{3n/2}+(3n/8)^{1/4}
 <\frac78T+\frac34\sqrt T<T.
\]
The asymptotic constant follows from $n\sim B^2/2$ and $M\sim B^3/6$.
\end{proof}

The elliptic core improves a leading constant, not the exponent or the
source separation: the rate still tends to zero and $T-A=1$.
Its number-field specialization guarantees infinitely many suitable
primes for each $\ell$, rather than every sufficiently large prime,
and gives no comparable effective field-size bound. The use of
algebraic torsion points avoids any claim of unbounded rational torsion.
